\section{Transport Layer (TCP \& UDP)}

\subsection{TCP (Transmission Control Protocol)}
\begin{itemize}
    \item Reliable, in-order byte stream; receiver ACKs received data.
    \item \textbf{vs. UDP:} UDP provides port+IP multiplexing (best effort) with no reliability, ordering, or flow control.
    \item \textbf{Sequence \#:} The byte number of the \textit{first} byte of data in the segment.
    \item \textbf{ACK \#:} The \textit{next expected} byte (Cumulative ACK).
    \item \textbf{MTU vs. MSS:}
    \begin{itemize}
        \item \textbf{MTU:} Max Transfer Unit (max packet size on link, e.g., 1500B).
        \item \textbf{MSS:} Max Segment Size (max payload). $MSS = MTU - (IP_{Header} + TCP_{Header})$. \\ Typically $MSS = MTU - 40$ (so $MSS < MTU$).
    \end{itemize}
	\item \textbf{TCP header (min 20B, + options):} src/dst ports, Seq \#, ACK \#, flags (SYN/ACK/FIN/RST/PSH/URG), window, checksum, urgent pointer.
\end{itemize}

\subsubsection{Connection Setup (3-Way Handshake)}
\begin{enumerate}
	\item Client $\to$ Server: \textbf{SYN}, Seq $= x$.\\
	\quad Announces client ISN (Initial Sequence Number, $x$) and requests connection.
	\item Server $\to$ Client: \textbf{SYN + ACK}, Seq $= y$, ACK $= x+1$.\\
	\quad Accepts, advertises server ISN, and ACKs client ISN.
	\item Client $\to$ Server: \textbf{ACK}, Seq $= x+1$, ACK $= y+1$.\\
	\quad Final ACK; both sides synchronized.
    \item (Connection Established)
\end{enumerate}

\subsubsection{Connection Teardown}
\begin{enumerate}
    \item Client sends \textbf{FIN}.
    \item Server replies \textbf{ACK}. (Client enters FIN\_WAIT\_2).
    \item Server sends \textbf{FIN}.
    \item Client replies \textbf{ACK}, waits \textbf{TIME\_WAIT} (usually 2*MSL), then closes.
\end{enumerate}

\subsubsection{Reliable data transfer (sender)}
\begin{itemize}
	\item Maintain sequence number and buffer; segment and send data.
	\item Send while window allows: $lastByteSent - lastByteAcked \le RW$ (receiver window).
	\item Start timer if not running; on timeout, retransmit the oldest unACKed segment.
\end{itemize}

\subsubsection{Reliable data transfer (receiver)}
\begin{itemize}
	\item Deliver in-order data and send cumulative ACK.
	\item If out-of-order segment arrives, send ACK for last in-order byte (duplicate ACK).
\end{itemize}

\subsubsection{Retransmissions}
\begin{itemize}
	\item Retransmit on either \textbf{timeout} or \textbf{duplicate ACKs}.
	\item \textbf{Fast retransmit:} many duplicate ACKs imply loss before timeout.
\end{itemize}

\subsubsection{RTT estimation}
\begin{itemize}
	\item Goal: set timeout slightly above RTT to avoid premature retransmits but still react fast to loss.
	\item \textbf{SampleRTT:} time from sending a segment until its first ACK returns (ignore retransmitted samples).
	\item Now we can compute a smoothed RTT which is exponentially weighted moving average of recent samples:
	\[ \textbf{EstimatedRTT} \leftarrow (1-\alpha)\,\text{EstimatedRTT} + \alpha\,\text{SampleRTT}\]
	typical value for $\alpha$ is $1/8$ (higher $\alpha$ reacts faster; lower is smoother).
	\item RTT variation (we add some safety margin):
	\[ \textbf{DevRTT} \leftarrow (1-\beta)\,\text{DevRTT} + \beta\,|\text{SampleRTT} - \text{EstimatedRTT}|\]
	typical value for $\beta$ is $1/4$ (higher $\beta$ reacts faster; lower is smoother).
	\item We may then set the timeout interval as:
	\[ 	\textbf{TimeoutInterval} = \text{EstimatedRTT} + 4\cdot \text{DevRTT} \]
\end{itemize}

\subsubsection{Flow Control vs. Congestion Control}
\begin{itemize}
	\item \textbf{Flow control}: prevents sender from overwhelming the receiver (receiver buffer limits).
	\item \textbf{Congestion control}: prevents sender from overloading the network.
	\item Sender can send while: $lastByteSent - lastByteAcked \le \min(RW, cwnd)$.
	\item When receiver buffer is the limiter, flow control dominates; when network is the limiter, congestion control dominates.
\end{itemize}

\subsubsection{Congestion Control Phases (AIMD)}
\textbf{AIMD (Additive Increase, Multiplicative Decrease):} Ensures stability and fairness.
\begin{enumerate}
    \item \textbf{Slow Start (SS):} Exponential growth; $cwnd = cwnd + 1$ per ACK (doubles every RTT).
    \item \textbf{Congestion Avoidance (CA):} Linear growth; $cwnd = cwnd + 1/cwnd$ per ACK ($\approx +1$ MSS per RTT).
\end{enumerate}
\begin{itemize}
    \item \textbf{Transition:} Threshold \textbf{ssthresh} decides SS vs CA: \\ If $cwnd < ssthresh \to$ SS. If $cwnd \ge ssthresh \to$ CA.
\end{itemize}

\subsubsection{TCP Tahoe}
\begin{itemize}
    \item On \textbf{any} loss (Timeout OR 3-Dup ACKs):
    \item Set $ssthresh = \lfloor cwnd/2 \rfloor$.
    \item Set $cwnd = 1$.
    \item Enter \textbf{Slow Start}.
\end{itemize}

\subsubsection{TCP Reno (Standard)}
\begin{itemize}
	\item Improves Tahoe for 3 duplicate ACKs (fast retransmit/fast recovery).
    \item On \textbf{Timeout:} Same as Tahoe ($cwnd=1$, restart SS etc).
    \item On \textbf{3 Duplicate ACKs (Fast Retransmit/Recovery):}
    \begin{enumerate}
        \item Retransmit missing segment immediately.
        \item Set $ssthresh = \lfloor cwnd/2 \rfloor$.
        \item Set $cwnd = ssthresh + 3$ (accounts for 3 packets buffering).
        \item Enter \textbf{Congestion Avoidance} (skips Slow Start).
    \end{enumerate}
	\item \textbf{Utilization intuition}
	\begin{itemize}
		\item Reno provides sawtooth behavior in $cwnd$ over time.
		\item With window size $W$ and RTT $RTT$, throughput is about $W/RTT$.
		\item If window halves after loss, average throughput over a cycle is roughly $0.75\,W/RTT$.
	\end{itemize}
\end{itemize}

\subsubsection{TCP Vegas}
\begin{itemize}
	\item Delay-based congestion control: compares expected throughput (from $cwnd/baseRTT$) to actual throughput (from $cwnd/RTT$).
	\item Adjusts $cwnd$ to keep the difference within a target range (\textbf{$\alpha$}, \textbf{$\beta$}), aiming to prevent queue buildup before loss.
	\item Less widely used because it competes poorly with loss-based flows (Reno), is sensitive to RTT measurement noise, and offers limited incremental deployability in the public Internet.
\end{itemize}

\subsection{UDP (User Datagram Protocol)}
\begin{itemize}
	\item Connectionless; no setup, no reliability, no ordering.
	\item Uses IP+port for multiplexing/demultiplexing (socket identified by destination port).
	\item Very small header; no congestion control, so can send as fast as the application wants.
	\item \textbf{UDP header (8B):} src/dst ports, length, checksum.
	\item Useful when timeliness matters more than reliability (e.g., streaming, VoIP), or when simplicity/overhead is key.
	\item Also used for request/response protocols like DNS and for DHCP (no TCP connection possible before address assignment).
	\item Reliability (if needed) is handled at the application layer.
\end{itemize}

\subsection{Reliable Data Transfer Protocols}
These protocols are the conceptual building blocks for reliability. While TCP is a specific protocol, it uses the mechanisms defined in these models.

\subsubsection{Stop-and-Wait}
\begin{itemize}
	\item \textbf{Mechanism:} Sender sends one packet, waits for ACK; if ACK/packet lost, timeout triggers retransmission.
    \item \textbf{Efficiency:} Very low in "long fat pipes" (high RTT). 
    \item \textbf{Utilization:} $U_{sender} = \frac{L/R}{RTT + L/R}$ \\ (where $L$=packet length, $R$=transmission rate).
	\item \textbf{Timeout lower bound:} $T_{out} \ge 2T_{prop} + T_{ack} + T_{pt}$.
	\item $T_{prop} = \frac{distance}{propagation\ speed}$, $T_{ack} = \frac{ACK\ size}{rate}$, \\ $T_{pt}$ = processing time.
	\item \textbf{Goodput (loss prob. $p$ for both data and ACK):} $\text{Goodput} = \frac{T_{packet}}{\frac{1}{(1-p)^2}(T_{packet}+T_{out})} = \frac{T_{packet}(1-p)^2}{T_{packet}+T_{out}}$.
	\item[] \textit{Note: Here we assume independent losses for data and ACK packets. If ACK losses are negligible, replace $(1-p)^2$ with $(1-p)$ in the Goodput formula, since now it's Geometric with success prob. $(1-p)$. \\ also, if $\overline{T_{out}}$ of successful transmissions is differenet than that of retransmissions, we would need to weight accordingly.}
	\item[] $\text{Goodput} = \frac{T_{packet}}{\left(\frac{1}{(1-p)^2}-1\right)(T_{packet}+T_{out})+(T_{packet}+\overline{T_{out}})}$
\end{itemize}

\textbf{Pipelined Protocols (Sliding Window):} \\
Allows multiple unacknowledged packets to be "in-flight" to increase utilization.
\subsubsection{Go-Back-N (GBN):}
\begin{itemize}
	\item \textbf{Sender:} Can send up to $N$ packets without ACKs. Uses a \textit{single timer} for the oldest unACKed packet.
	\item \textbf{Receiver:} Uses \textbf{Cumulative ACKs}. If a packet is missing, it discards all subsequent out-of-order packets.
	\item \textbf{Receiver window:} size 1; accepts only next in-order packet.
	\item \textbf{On Timeout:} Sender retransmits \textit{all} $N$ transmitted but unACKed packets.
	\item \textbf{Goodput (loss prob. $p$):} let $a=\frac{T_{out}+T_{packet}}{T_{packet}}$, then $\text{Goodput}_{GBN} = \frac{1-p}{1-p+p a}$.
\end{itemize}
\subsubsection{Selective Repeat (SR):}
\begin{itemize}
	\item \textbf{Sender:} Maintains a \textit{timer for each} individual packet.
	\item \textbf{Receiver:} ACKs packets individually. \textbf{Buffers} out-of-order packets for later delivery.
	\item \textbf{On Timeout:} Sender retransmits \textit{only} the specific lost packet.
	\item \textbf{Window Size Constraint:} To avoid sequence ambiguity, $WindowSize \le \frac{1}{2} \cdot MaxSeqNum$.
\end{itemize}

\begin{center}
\renewcommand{\arraystretch}{1.15}
\setlength{\tabcolsep}{3pt}
\begin{tabular}{p{0.22\columnwidth}|p{0.34\columnwidth}|p{0.34\columnwidth}}
\textbf{Criteria} & \textbf{Selective Repeat} & \textbf{Go-Back-N} \\
\hline
Bandwidth utilization & Retransmits only lost packets & May retransmit already received packets \\
Window sizes & Sender: $N$, Receiver: $N$ & Sender: $N$, Receiver: $1$ \\
Sequence numbers & $2N$ & $N+1$ \\
Out-of-order handling & Buffers and tracks out-of-order packets & No buffering; accepts only in-order \\
Receiver storage & Stores out-of-order until missing arrives & No need to store (sender may resend) \\
\end{tabular}
\end{center}

\subsubsection{Implementation \& Relation to TCP}
\begin{itemize}
    \item \textbf{Where they run:} These protocols run in the \textbf{Transport Layer} (Layer 4) within the \textbf{Operating System} of the end-hosts. Routers and switches (the "network core") generally do not implement these.
    \item \textbf{TCP Implementation:} 
    \begin{itemize}
        \item TCP is a hybrid: It uses \textbf{Cumulative ACKs} (like GBN) but typically \textbf{Buffers out-of-order segments} and can use \textbf{SACK} (Selective ACK) to retransmit only missing data (like SR).
    \end{itemize}
\end{itemize}